We evaluated an open 8B LLM for clinical note generation on the official ACI-Bench and MTS-Dialog test sets
and compared it with published MEDIQA-Chat 2023 systems using the official ROUGE configuration. Untrained
Llama-3.1-8B-Instruct matched GPT-4 zero-shot on both datasets (ROUGE-1 0.519 and 0.303). QLoRA fine-tuning
raised ROUGE on the short MTS-Dialog sections to the level of fine-tuned shared-task systems (0.400), but
increased the entity-level hallucination rate on both datasets (ACI-Bench 14.4\%$\rightarrow$51.7\%,
MTS-Dialog 20.4\%$\rightarrow$35.6\%); adding the fact table to the fine-tuned model did not bring it back
to the level of the untrained model. FaCTS addresses this without training: conditioning the model on a
negation- and dose-aware fact table lowered hallucination to 10.6\% and 15.8\% with no loss of coverage, giving
the best fact-level F1 on both datasets, and sentence-level NLI verification lowered it further to 8.7\% and
14.2\% at a cost of about 11 points of coverage.

Future work follows directly from these results. (1) \emph{Repair instead of drop}: rewrite red sentences
with the generator, as in detector-guided self-refinement~\cite{seethakantha2026hdsr}, to keep their
supported facts. (2) \emph{Factuality preference training}: train an ORPO~\cite{hong2024orpo} adapter on pairs
of reference notes and copies with one injected error (negation flip, dose change, entity swap, inserted
fact), so that the model, not only the pipeline, learns to prefer grounded notes. (3) A better verifier for
dialogue negation, e.g.\ fine-tuning the NLI model on clinical question--answer pairs. (4) Clinician
evaluation of a sample of notes, and the shared-task BERTScore and BLEURT for a complete comparison.
(5) Rebalancing the fine-tuning data towards full notes and selecting checkpoints by factuality rather than
validation loss.
